\documentclass[a4paper,11pt]{article}
\usepackage{lipsum}
\pagestyle{empty}
\setlength{\parindent}{0pt}
\setlength{\parskip}{0pt}

% Minimal test: ONE figure, ONE long paragraph that WILL break across pages
% This should trigger ghost_widen for the continuation

\directlua{
  local widen_log = {}
  local orig_postlb = luatexbase.remove_from_callback("post_linebreak_filter", "swarmwrap: penalty at parshape boundary")
  function swarmwrap_post_lb(head, groupcode)
    if swarmwrap_ghost_widen then
      local lw_val = swarmwrap_safe_dimen("swarmwrap@lw@saved")
      local ghost_tw = swarmwrap_ghost_tw
      local tol = 15.0
      local narrow_max = (ghost_tw / 65536.0) + tol
      local lines_found = 0
      local lines_widened = 0
      local widths = {}
      local current = head
      while current do
        if current.id == 0 then
          lines_found = lines_found + 1
          local lw = current.width / 65536.0
          widths[#widths + 1] = string.format("%.1f", lw)
          if lw <= narrow_max and lw > 0 then
            current.width = lw_val
            local inner_head = current.head
            if inner_head then
              local inner_last = node.tail(inner_head)
              local fil = node.new(node.id("glue"))
              fil.stretch = 65536
              fil.stretch_order = 2
              if inner_last then
                node.insert_after(inner_head, inner_last, fil)
              end
            end
            lines_widened = lines_widened + 1
          end
        end
        current = current.next
      end
      widen_log[#widen_log + 1] = "GW: found=" .. lines_found .. " widened=" .. lines_widened .. " w={" .. table.concat(widths, ",") .. "} gtw=" .. string.format("%.1f", ghost_tw / 65536.0)
      swarmwrap_ghost_widen = false
      return head
    end
    return orig_postlb(head, groupcode)
  end
  luatexbase.add_to_callback("post_linebreak_filter",
    function(head, groupcode)
      local ok, result = pcall(swarmwrap_post_lb, head, groupcode)
      if ok then return result end
      return head
    end, "swarmwrap: penalty at parshape boundary")
}

\newenvironment{swarmwrap}{%
  \begin{lrbox}{\swarmwrap@box}%
}{%
  \end{lrbox}%
}
\newcommand{\swarmwrapnext}{}

\begin{document}
\vspace*{5cm}
\noindent\textbf{Start here.}\par
% Manually set up a parshape for 10 narrow lines, then full width
\parshape 10 0pt 220pt 0pt 220pt 0pt 220pt 0pt 220pt 0pt 220pt 0pt 220pt 0pt 220pt 0pt 220pt 0pt 220pt 0pt 220pt 0pt 360pt
This is line 1 of narrow text that should trigger ghost widening when it breaks across pages.
This is line 2 of narrow text continuing the paragraph before it breaks.
This is line 3 of narrow text, getting closer to the page break.
This is line 4 of narrow text, almost there.
This is line 5 of narrow text, should be near the page break.
This is line 6 of narrow text, right at the page break boundary.
This is line 7 of narrow text, might split across pages.
This is line 8 of narrow text, continuation lines start here.
This is line 9 of narrow text, more continuation lines.
This is line 10 of narrow text, last narrow line.
This is line 11, back to full width after the narrow zone.
This is line 12, also full width, should not be affected.
This is line 13, full width, normal text continues here.
This is line 14, more full width text on this page.
\parshape 1 0pt 360pt
This is a new paragraph starting at full width. Should not be affected.
More full-width text here to fill the page.
More full-width text here to fill the page.
More full-width text here to fill the page.
\end{document}\parshape 10 0pt 220pt 0pt 220pt 0pt 220pt 0pt 220pt 0pt 220pt 0pt 220pt 0pt 220pt 0pt 220pt 0pt 220pt 0pt 220pt 0pt 360pt
\lipsum[1-20]%
